\documentclass[11pt]{article}
\usepackage[a4paper,margin=1in]{geometry}
\usepackage{amsmath,amssymb,amsthm}
\usepackage[T1]{fontenc}
\usepackage{lmodern}
\usepackage[hidelinks]{hyperref}
\setlength{\emergencystretch}{3em}

\theoremstyle{plain}
\newtheorem{theorem}{Theorem}
\newtheorem{lemma}[theorem]{Lemma}
\newtheorem{proposition}[theorem]{Proposition}
\newtheorem{corollary}[theorem]{Corollary}
\theoremstyle{definition}
\newtheorem{definition}[theorem]{Definition}
\theoremstyle{remark}
\newtheorem{remark}[theorem]{Remark}

\title{A Counterexample to Conjecture 00000002040:\\
Largest Sum-Free Subsets of $[n]$ Have Density $\to 1/2$, Not $1/4$}
\author{%
  Feiyu\thanks{Independent researcher.}
  \and
  Claude (Opus 5.5)\thanks{Anthropic.}%
}
\date{2026-10-02}

\begin{document}
\maketitle

\begin{abstract}
Conjecture 00000002040 asserts that the density of the largest 3-independent
(sum-free) subset of $[n]=\{1,\ldots,n\}$ tends to a limit that is at most, and
in fact exactly, $1/4$. The largest such subset has exactly
$\lceil n/2\rceil$ elements, so the density tends to $1/2$. The odd numbers
show that $1/2$ is attained, and a short reflection argument shows it cannot be
beaten. Everything is formalized in Lean~4 against Mathlib.
\end{abstract}

\section{The conjecture as stated}

The original text of conjecture 00000002040 reads:

\begin{quote}
Definition: A 3-independent set has no solutions of $x + y = z$. Conjecture:
The density limit of maximal 3-independent sets in $[n]$ exists and is at most
$1/4$ (with correction $0$); the limit is exactly $1/4$ (attained by the
natural ternary cascade construction).
\end{quote}

\section{Reading of the statement}

A set $A$ of positive integers is \emph{3-independent} if there are no
$x,y,z\in A$ (not necessarily distinct) with $x+y=z$; such sets are usually
called sum-free. Let
\[
  M(n)=\max\{|A| : A\subseteq[n]\text{ 3-independent}\},
  \qquad d(n)=M(n)/n .
\]
``Maximal'' refers to the largest such sets; the Chinese version of the
statement (in the conjecture file) says so explicitly, speaking of the largest
3-independent set in $[n]$. The conjecture asserts that $d(n)$
converges, that its limit is at most $1/4$, and that the limit is exactly
$1/4$. The parenthetical remarks (``with correction~$0$'', ``attained by the
natural ternary cascade construction'') add no mathematical content that we
need.

\section{Result}

\begin{theorem}\label{thm:main}
For every $n\ge0$, $M(n)=\lceil n/2\rceil=\lfloor (n+1)/2\rfloor$. Consequently
$d(n)\to\frac12$, so $d(n)$ has no limit $\le\frac14$, and Conjecture 00000002040
is false.
\end{theorem}

\section{Proof}

\begin{lemma}[lower bound]\label{lem:odd}
The set of odd numbers in $[n]$ is 3-independent and has
$\lfloor(n+1)/2\rfloor$ elements.
\end{lemma}

\begin{proof}
A sum of two odd numbers is even, hence not odd. The odd numbers in $[n]$ are
$2k+1$ for $0\le k<\lfloor(n+1)/2\rfloor$.
\end{proof}

\begin{lemma}[upper bound]\label{lem:upper}
Every 3-independent $A\subseteq[n]$ has $|A|\le\lfloor(n+1)/2\rfloor$.
\end{lemma}

\begin{proof}
We may assume $A\ne\emptyset$; let $m=\max A$. For $a\in A$ with $a<m$, the
number $m-a$ is not in $A$, since otherwise $a+(m-a)=m$ would be a solution
inside $A$. So $A$ and $B=\{m-a : a\in A,\ a<m\}$ are disjoint. Both lie in
$[m]$, and $|B|=|A|-1$. Hence $2|A|-1\le m\le n$.
\end{proof}

\begin{proof}[Proof of Theorem~\ref{thm:main}]
Lemmas~\ref{lem:odd} and~\ref{lem:upper} give $M(n)=\lfloor (n+1)/2\rfloor$.
For $n\ge1$ this gives $\frac12\le d(n)\le\frac12+\frac1{2n}$, so $d(n)\to\frac12$. A
convergent sequence has a unique limit, and $\frac12>\frac14$.
\end{proof}

\section{Formalization}

The accompanying Lean~4 project (\texttt{lean/Submission/Basic.lean}) states
the conjecture and proves its negation against Mathlib.

\begin{itemize}
  \item \texttt{IsThreeIndependent A} --- no $x,y,z\in A$ with $x+y=z$;
    \texttt{maxSize n} --- the largest size of a 3-independent subset of
    \texttt{Finset.Icc 1 n}, a supremum over its powerset;
    \texttt{density n} $=\texttt{maxSize n}/n$ in $\mathbb R$.
  \item \texttt{ConjectureHolds} --- there is $L\le\frac14$ with
    \texttt{density} $\to L$, and \texttt{density} $\to\frac14$.
  \item \texttt{odds}, \texttt{card\_odds}, \texttt{isThreeIndependent\_odds}
    --- Lemma~\ref{lem:odd}.
  \item \texttt{card\_le\_of\_isThreeIndependent} --- Lemma~\ref{lem:upper};
    \texttt{maxSize\_eq} --- $M(n)=\lfloor(n+1)/2\rfloor$.
  \item \texttt{half\_le\_density}, \texttt{density\_le},
    \texttt{tendsto\_density} --- $d(n)\to\frac12$.
  \item \texttt{conjecture\_00000002040\_false} ---
    $\lnot\,$\texttt{ConjectureHolds}, by uniqueness of limits.
\end{itemize}

The toolchain is \texttt{leanprover/lean4:v4.33.1}, Mathlib is pinned to
\texttt{v4.33.1}, and \texttt{lake build} succeeds. The project contains no
\texttt{sorry}, no \texttt{admit} and no \texttt{native\_decide}, and
introduces no axioms:
\begin{center}
\texttt{\#print axioms conjecture\_00000002040\_false}
\end{center}
reports exactly \texttt{propext}, \texttt{Classical.choice} and
\texttt{Quot.sound}.

\section{Remarks}

The value $\lceil n/2\rceil$ is classical. Besides the odd numbers, the upper
half $\{\lfloor n/2\rfloor+1,\ldots,n\}$ is also extremal. The constant $1/4$
does not arise for sum-free subsets of $[n]$. For comparison, in the cyclic
group $\mathbb Z/p\mathbb Z$ ($p$ prime) the largest sum-free subsets have
$\lfloor(p+1)/3\rfloor$ elements, density about $1/3$.

\end{document}
